% =====================================================================
% Subsecao 5.5 — Avaliacao Empirica (Rodada 2)
%
% Como usar:
%   1. Substitua todos os <PREENCHER: ...> pelos numeros reais dos CSVs.
%   2. Copie matriz_confusao.png para images/ do projeto LaTeX.
%   3. Inclua este arquivo no tcc_mikhael.tex apos a Subsecao 5.4 com:
%        \input{secoes/secao_5_5_empirica}
%      (ou cole o conteudo direto).
%   4. Atualize a Subsecao 6.2 para citar os numeros desta secao
%      (use \ref{tab:h1-percentis}, \ref{tab:h2-disponibilidade},
%       \ref{tab:classificador-global} e \ref{fig:matriz-confusao}).
% =====================================================================

\subsection{Avaliacao Empirica das Hipoteses H1 e H2}
\label{sec:avaliacao-empirica}

Esta subsecao apresenta os resultados quantitativos obtidos com a Rodada~2
de benchmarks, conduzida em ambiente local com o \texttt{docker compose}
do Fyna (servicos \texttt{fyna-back} e \texttt{fyna-ai}). Os artefatos de
execucao estao em \texttt{docs/benchmarks/} e os dados brutos em
\texttt{resultados\_h1\_*.csv}, \texttt{resultados\_h2\_*.csv} e
\texttt{resultados\_brutos.csv}.

\paragraph{Ambiente de teste.}
\begin{itemize}
  \item Hardware: <PREENCHER: CPU/RAM/Disco do host>
  \item SO: <PREENCHER: ex.: Windows 11 / Linux Ubuntu 22.04>
  \item Containers: \texttt{fyna-backend} (Spring Boot 3.x), \texttt{fyna-ai-engine} (FastAPI), Postgres no host
  \item Carga: Locust \texttt{<PREENCHER: versao>}, 25 usuarios virtuais, ramp-up 5 VUs/s, 3 minutos
  \item Data da coleta: <PREENCHER: AAAA-MM-DD>
\end{itemize}

% =====================================================================
\subsubsection{H1 — Latencia (p95 < 200 ms)}
\label{sec:h1}

A Tabela~\ref{tab:h1-percentis} sintetiza os percentis de latencia por
endpoint, extraidos de \texttt{resultados\_h1\_stats.csv}.

\begin{table}[!htb]
  \centering
  \caption{Percentis de latencia (ms) — 25 VUs, 3 min.}
  \label{tab:h1-percentis}
  \begin{tabular}{lrrrrr}
    \hline
    \textbf{Endpoint} & \textbf{p50} & \textbf{p95} & \textbf{p99} & \textbf{Req/s} & \textbf{Erros (\%)} \\
    \hline
    POST /api/v1/transactions    & <PREENCHER> & <PREENCHER> & <PREENCHER> & <PREENCHER> & <PREENCHER> \\
    GET  /api/v1/accounts        & <PREENCHER> & <PREENCHER> & <PREENCHER> & <PREENCHER> & <PREENCHER> \\
    GET  /api/v1/transactions/\{id\} & <PREENCHER> & <PREENCHER> & <PREENCHER> & <PREENCHER> & <PREENCHER> \\
    \hline
    \textbf{Agregado}            & <PREENCHER> & <PREENCHER> & <PREENCHER> & <PREENCHER> & <PREENCHER> \\
    \hline
  \end{tabular}
\end{table}

\paragraph{Veredito H1.}
<PREENCHER: APROVADO ou REPROVADO — escreva o p95 agregado e compare com 200~ms>.
<PREENCHER: comentario curto sobre o endpoint mais lento, possiveis gargalos
(p.\,ex., serializacao, indices, latencia do Postgres) e qualquer mitigacao
ja adotada na implementacao>.

% =====================================================================
\subsubsection{H2 — Disponibilidade sob falha do microsservico de IA}
\label{sec:h2}

Para validar H2 foi injetada uma falha controlada parando o servico
\texttt{fyna-ai} via \texttt{docker compose stop} por 120~s, com fases
baseline e recuperacao de 60~s cada. Durante todo o periodo o backend
recebeu requisicoes \texttt{POST /api/v1/transactions} sem
\texttt{categoryId} (caminho que dispara classificacao \texttt{@Async} no
\texttt{AIEngineClient}, sujeita ao circuit~breaker). A
Tabela~\ref{tab:h2-disponibilidade} resume os resultados.

\begin{table}[!htb]
  \centering
  \caption{Taxa de sucesso por fase do experimento H2.}
  \label{tab:h2-disponibilidade}
  \begin{tabular}{lrrr}
    \hline
    \textbf{Fase} & \textbf{Total req.} & \textbf{Sucesso (2xx)} & \textbf{Taxa (\%)} \\
    \hline
    Baseline (60 s, IA UP)             & <PREENCHER> & <PREENCHER> & <PREENCHER> \\
    Falha injetada (120 s, IA DOWN)    & <PREENCHER> & <PREENCHER> & <PREENCHER> \\
    Recuperacao (60 s, IA UP)          & <PREENCHER> & <PREENCHER> & <PREENCHER> \\
    \hline
  \end{tabular}
\end{table}

\paragraph{Veredito H2.}
<PREENCHER: APROVADO ou REPROVADO — escreva a taxa de sucesso da fase de
falha e compare com o limiar de 99\%>.
<PREENCHER: comente sobre o impacto observado em
\texttt{AIEngineClient.classifyTransaction} (chamada \texttt{@Async} + circuit~breaker:
5 falhas / 60~s) e o fato de a transacao ser persistida antes do trigger
do classificador, garantindo que falhas no \texttt{fyna-ai} nao propaguem
para o usuario final>.

% =====================================================================
\subsubsection{Avaliacao do classificador supervisionado}
\label{sec:classificador}

Foram geradas \texttt{<PREENCHER: N>} transacoes sinteticas em PT-BR,
distribuidas entre as categorias canonicas do Fyna. Cada transacao foi
criada via \texttt{POST /api/v1/transactions} sem \texttt{categoryId} e a
predicao do classificador foi coletada em
\texttt{GET /api/v1/ai/classifications/transaction/\{id\}}.

\begin{table}[!htb]
  \centering
  \caption{Metricas globais do classificador.}
  \label{tab:classificador-global}
  \begin{tabular}{lr}
    \hline
    \textbf{Metrica}        & \textbf{Valor} \\
    \hline
    Acuracia global         & <PREENCHER> \\
    F1 (macro)              & <PREENCHER> \\
    Precisao (macro)        & <PREENCHER> \\
    Revocacao (macro)       & <PREENCHER> \\
    Latencia p50 (ms)       & <PREENCHER> \\
    Latencia p95 (ms)       & <PREENCHER> \\
    Versao do modelo        & <PREENCHER> \\
    \hline
  \end{tabular}
\end{table}

\begin{table}[!htb]
  \centering
  \caption{Metricas por categoria (extraido de
  \texttt{metricas\_por\_categoria.csv}).}
  \label{tab:classificador-por-categoria}
  \begin{tabular}{lrrrr}
    \hline
    \textbf{Categoria} & \textbf{Precision} & \textbf{Recall} & \textbf{F1} & \textbf{Support} \\
    \hline
    Alimentação    & <PREENCHER> & <PREENCHER> & <PREENCHER> & <PREENCHER> \\
    Transporte     & <PREENCHER> & <PREENCHER> & <PREENCHER> & <PREENCHER> \\
    Mercado        & <PREENCHER> & <PREENCHER> & <PREENCHER> & <PREENCHER> \\
    Moradia        & <PREENCHER> & <PREENCHER> & <PREENCHER> & <PREENCHER> \\
    Saúde          & <PREENCHER> & <PREENCHER> & <PREENCHER> & <PREENCHER> \\
    Lazer          & <PREENCHER> & <PREENCHER> & <PREENCHER> & <PREENCHER> \\
    Educação       & <PREENCHER> & <PREENCHER> & <PREENCHER> & <PREENCHER> \\
    Salário        & <PREENCHER> & <PREENCHER> & <PREENCHER> & <PREENCHER> \\
    Investimentos  & <PREENCHER> & <PREENCHER> & <PREENCHER> & <PREENCHER> \\
    \hline
  \end{tabular}
\end{table}

\begin{figure}[!htb]
  \centering
  \includegraphics[width=0.85\linewidth]{images/matriz_confusao.png}
  \caption{Matriz de confusao normalizada por linha. A diagonal indica a
  fracao de instancias da categoria real classificadas corretamente.}
  \label{fig:matriz-confusao}
\end{figure}

\paragraph{Discussao.}
<PREENCHER: discuta categorias com maior confusao mutua (p.\,ex.,
"Alimentacao" vs "Mercado"), efeito do limiar de confianca configurado
em \texttt{settings.classifier\_confidence\_threshold} e o impacto do
aprendizado por correcoes (\texttt{was\_corrected=true}) descrito em
\texttt{classifier.py}>.

% =====================================================================
\subsubsection{Sintese e implicacoes para H1 e H2}
\label{sec:sintese-empirica}

<PREENCHER: 1-2 paragrafos relacionando os resultados acima com as
proposicoes do Capitulo~6. Indique se as evidencias coletadas sustentam
ou nao H1 (latencia) e H2 (disponibilidade) e referencie as tabelas
deste capitulo nas Subsecoes 6.2.1 e 6.2.2.>
